\textbf{\textit{Soal. }}In the right triangle $ABC,$ $CF$ is the altitude based on the hypotenuse $AB.$
The circle centered at $B$ and passing through $F$ and the circle with centre $A$ and the same radius intersect at a point of $CB.$
Determine the ratio $FB : BC.$

\begin{proof}[\textbf{Solusi. }] Misalkan $\angle CBA = \alpha$ dan $BD=BF=DA=x$. Karena $BD=DA$ kita punya $\angle DAB = \alpha$ dan $\angle CDA = 2\alpha$. Ini berarti $CA = DA \sin \angle CDA = x \sin 2\alpha$. Lalu, karena $CF \perp AB$ maka $\angle BCF = 90^\circ-\alpha$ dan $\angle FCA=\alpha$. Oleh karena itu, $FA = CA \sin \angle FCA = x\sin 2\alpha \sin \alpha$.
    \begin{center}
    \begin{asy}
        import olympiad;
        import geometry;
        size(150);

        pair A,B,C,D,F,M;
        C = (0,0);
        B = (0,4);
        A = (3.0619678, 0);
        F = foot(C, A, B);
        M = (A+B)/2;
        line l = perpendicular(M, line(B,A));
        D = intersectionpoint(line(B,C),l);
        dot("$A$", A, SE);
        dot("$B$", B, N);
        dot("$C$", C, SW);
        dot("$F$", F, NE);
        dot("$D$", D, W);
        draw(A--B--C--A);
        draw(B--D, red);
        draw(D--A, red);
        draw(B--F, red);
        draw(C--F);
        draw(rightanglemark(C,F,A,4), blue);
    \end{asy}
\end{center}

Di lain sisi, karena $BD=DA$, kita juga punya $\frac{1}{2}AB = BD \cos \angle DBA = x \cos \alpha$. Berarti kita dapatkan
\begin{align*}
    AB &= BF + FA\\
    2x \cos \alpha &= x + x\sin 2\alpha \sin \alpha\\
    2 \cos \alpha &= 1 + \sin 2\alpha \sin \alpha\\
    2 \cos \alpha &= 1 + (2\sin \alpha \cos \alpha \sin \alpha)\\
    2 \cos \alpha &= 1 + (2\sin^2 \alpha \cos \alpha)\\
    2 \cos \alpha &= 1 + (2(1 - \cos^2 \alpha) \cos \alpha)\\
    2 \cos \alpha - (2(1 - \cos^2 \alpha) \cos \alpha) &= 1 \\
    2( 1 - (1 - \cos^2 \alpha) \cos \alpha) &= 1 \\
    2\cos^3 \alpha &= 1\\
    \cos \alpha &= \dfrac{1}{\sqrt[3]{2}}.
\end{align*}
Sekarang, perhatikan bahwa $CD = DA \cos \angle CDA = x \cos 2\alpha$. Berarti $BC=BD+DC=x+x\cos2\alpha=x(1+\cos2\alpha)=x(1+2\cos^2\alpha-1)=2x\cos^2\alpha$ sehingga kita punya
\begin{align*}
    \dfrac{FB}{BC} = \dfrac{x}{2x\cos^2\alpha}=\dfrac{1}{2\cdot\left(\frac{1}{\sqrt[3]{2}}\right)^2}=\dfrac{1}{\sqrt[3]{2}}.
\end{align*}
\end{proof}